\documentclass[../../main.tex]{subfiles}
\begin{document}
\chapter{Groups and homomorphisms}
\label{ch:groups}

\begin{objectives}
\begin{itemize}
  \item Recognize the group axioms.
  \item Use cancellation in a proof.
  \item Describe a homomorphism and its kernel.
\end{itemize}
\end{objectives}

\section{The group axioms}
The following examples illustrate how to write a sequence of definitions,
results, and proofs. For a fuller treatment, see \textcite{judson2020}.
\sidenote{A definition specifies the assumptions available in every later
proof. Keep track of which axiom justifies each step.}

\begin{definition}[Group]\label{def:group}
A \vocab{group} is a set $G$ with a binary operation $G\times G\to G$
that is associative, has an identity $e$, and gives every $g\in G$ an
inverse $g\inv$. A group is \vocab{abelian} if $gh=hg$ for all $g,h\in G$.
\end{definition}

\begin{eg}
The integers $\ZZ$ form an abelian group under addition. The identity is
$0$, and the inverse of $n$ is $-n$. The nonzero real numbers form a group
under multiplication, but all real numbers do not: $0$ has no inverse.
\end{eg}

\begin{lemma}[Cancellation]\label{lem:cancellation}
If $a,b,c\in G$ and $ab=ac$, then $b=c$.
\end{lemma}
\begin{proof}
Multiply both sides on the left by $a\inv$ and use associativity:
\[
\begin{aligned}
  b&=(a\inv a)b=a\inv(ab)\\
   &=a\inv(ac)=(a\inv a)c=c.
\end{aligned}
\]
\end{proof}

\begin{bigidea}[Undo an operation]\label{idea:cancellation}
An inverse lets us undo multiplication on the same side. Cancellation is
therefore a consequence of the group axioms, even when multiplication is
not commutative.
\end{bigidea}

\begin{theorem}[Uniqueness of inverses]\label{thm:unique-inverse}
Every element of a group has exactly one inverse.
\end{theorem}
\begin{proof}
Existence is part of \cref{def:group}. If $b$ and $c$ are inverses of $a$,
then $ab=e=ac$, so \cref{lem:cancellation} gives $b=c$.
\end{proof}

\begin{corollary}
For every $a\in G$, we have $(a\inv)\inv=a$.
\end{corollary}

\section{Maps between groups}
\begin{definition}[Homomorphism]\label{def:homomorphism}
A map $f:G\to H$ is a \vocab{homomorphism} if
$f(xy)=f(x)f(y)$ for all $x,y\in G$. Its kernel is
\[
  \Ker f\defeq\{g\in G:f(g)=e_H\}.
\]
The set of homomorphisms from $G$ to $H$ is denoted by $\Hom(G,H)$.
\end{definition}

\begin{prop}\label{prop:identity-map}
A homomorphism $f:G\to H$ sends the identity of $G$ to the identity of $H$.
\end{prop}
\begin{proof}
Since $f(e_G)=f(e_Ge_G)=f(e_G)f(e_G)$, cancellation in $H$ gives
$e_H=f(e_G)$.
\end{proof}

\begin{explanation}
For the reduction map $\pi:\ZZ\to\Zc{n}$, addition is preserved:
$\pi(a+b)=\pi(a)+\pi(b)$. Its kernel is $n\ZZ$. The notation
$\Zc{n}$ is produced by \verb|\Zc{n}| from \texttt{macro.sty}.
Here $n$ is a positive integer.
\end{explanation}

\begin{remark}
For a general group, a subgroup must be normal before a quotient group is
defined. The command \verb|\normalin| prints the normal-subgroup symbol:
$N\normalin G$.
\end{remark}

\section{Exercises}
\begin{problem}[Identity]
Show that the identity element of a group is unique.
\end{problem}
\begin{soln}
If $e$ and $e'$ are both identity elements, then $e=ee'=e'$.
\end{soln}

\begin{exercise}[Inverses of products]\label{ex:product-inverse}
Prove that $(ab)\inv=b\inv a\inv$. Check both products with $ab$.
\end{exercise}

\begin{recap}
The group axioms give cancellation and uniqueness of inverses.
Homomorphisms preserve the operation and the identity; their kernels
record which elements map to the identity. In \cref{ex:product-inverse},
the order of the inverse factors matters.
\end{recap}
\end{document}
